
% 来源：EPUB part049.xhtml
\section{不做空、不借钱、不懂不碰}

“市场就像上帝一样，帮助那些自己帮助自己的人，但与上帝不一样的地方是，他不会原谅那些不知道自己在做什么的人。”（巴菲特，1983）

% 来源：EPUB part050.xhtml
\subsection{不懂不碰是铁律}

老调重弹：不做空、不用margin，不做不懂的东西。什么时候明白都不晚，不过真的明白恐怕都是在付出大代价之后。（2012-03-14）

买自己不懂的东西是为什么85\%的人可以在牛市熊市都亏钱的原因。没有比因为股票涨了很多而买入更荒唐的理由了。（2015-06-26）

我投资中的不为清单（Stop Doing List）很简单，就是不懂不碰。懂主要集中在两个方面：商业模式和企业文化。（2020-10-10）

\BookQuestion{网友：好想买能安心睡觉的公司。}

能不能安心睡觉在买以前就决定了。（2011-03-29）

我只是侥幸看懂了几个我刚好懂得的好企业而已。绝大部分我也是看不懂的，和绝大多数人一样。不过我看不懂就不碰了，很多人做不到还想赚快钱。当然，我也是很想赚快钱的，但知其不可为而已。（2024-06-17）

不懂不做，巴菲特一直如此。但懂有时候要有个过程，巴菲特也不例外，也会碰到以为懂了但后来发现其实不完全懂的东西。另外，巴菲特公司还有其他人管投资，巴菲特是放权给他们做的，但他们的理解有时候和他会不完全一样，所以大家看到的巴菲特的持股未必一定就是巴菲特的风格。其实比亚迪就不是巴菲特的风格，通用汽车也不是。在石油公司上巴菲特也认为自己犯了一个几十亿美元的错误（把中石油上赚的钱亏进去很多）。（2012-10-10）

巴菲特讲的能力圈子指的就是每个人都有自己有限的了解范围。每个人总是有些东西懂有些东西不懂。所以投自己懂的东西就容易知道价值，就知道什么叫便宜，就有机会赚钱，反之亦然。如果你买觉得心虚的东西就是投机，投机当然也有赚钱的时候，但风险大，睡不好觉。（2010-03-25）

\BookQuestion{网友：如果一个企业回购自己的股票，是否就可以认为这家公司便宜了，是否该跟进买入？}

任何“跟进”都是错的，因为你早晚会错。

跟进的意思就是自己不明白瞎买的意思，所以早晚会跟错的。而且，跟进也很难赚钱，因为不是很明白，所以拿不住。如果你也明白就不一样了。（2010-07-23）

今天在网上看到有人说跟着我买阿里，说这话的人都是投机分子！自己不懂的东西千万别跟着别人投，拉人壮胆也是没用的哈，不管这人是谁。反过来也一样，自己懂的东西不怕别人不喜欢，哪怕那人是巴菲特或芒格也没关系。（2014-09-23）

股市上那些长期亏钱的大多属于不知道自己能力圈有多大的人。（2012-12-29）

知道自己能力圈有多大或者说知道自己能力圈的边界在哪里，要远远重要过能力圈有多大，这也是人们能看到很多很“聪明”的人投资表现长期不好的原因。当然，这些“聪明”人会把别人的成功或自己的不成功归结于运气或accidents（意外），而且他们总是能够很“聪明”地找到办法让自己认为确实如此。《功夫熊猫1》里说“There are no accidents（世事无偶然）”其实是蛮有道理的。（2013-02-01）

“投资并非一个智商为160的人就一定能击败智商为130的人的游戏。”（巴菲特）

大概智商高的人未必知道自己的能力圈边界在哪里，但智商高的人可能往往容易越出自己的能力圈。

“我们之所以取得目前的成就，是因为我们关心的是寻找那些我们可以跨越的一英尺栏，而不是去拥有什么能飞越七英尺的能力。”（巴菲特）

这说的也是能力圈的问题。经常看到人们在讨论本身一头雾水的投资概念，让我常常想起这话。（2012-07-28）

做过企业家的人本来可以有一个优势，那就是更容易看懂公司，更容易知道自己的能力圈在哪里。但我发现我认识的很多企业家对股市并没有任何认识，和普通人无异，同时大多数人是不碰股市的（因为他们认为自己不懂）。不懂不碰的就是好的投资者，哪怕一点都不碰也是。

仔细想来，做过企业家的人未必就更容易知道自己能力圈的大小，有可能会相反。（2011-09-07）

\BookQuestion{网友：投资中超出自己的领域时，求助专家建议也许不是最坏的打算？}

“砖家”建议你也敢听？超出领域时最好别碰，除非你能搞懂。当你不懂时，专家帮不了忙。专家只能提供参考意见，你要是不懂还是无法决策的。（2010-04-18）

专家可以在细节上帮很多忙，但不能帮你决策。（2010-04-19）

\BookQuestion{网友：您能谈谈亏损的经历吗？大家都很想知道您投资失败的案例。}

我买过一只航空股，叫FRNT，亏了很多钱（之前有过一只航空股赚过很多钱）。这家公司现金流很好，就是有点太小，结果2008年被信用卡公司突然要全额现金（因为有4家别的小航空公司突然破产）。当时如果CEO想办法给对方现金的话（当时可以卖些飞机和其他资产的），其实也能过关。结果这个CEO突然宣布进破产法庭，最后该给人的钱还是给了，把公司真搞破产了。当时买这家公司时，主要是觉得石油太贵了，一定会掉下来。石油一掉，航空公司最先受益。结果都对了，油也确实掉下来了，这家公司也确实后来赚了很多钱，但我们却亏掉了。

我这里犯的还是做了自己其实不明白的事。比较遗憾的是，之前巴菲特已经告诉我小心航空公司。这家公司破产后我真正想通了巴菲特是什么意思。所以以后不会再轻易碰航空公司了。（2010-03-27）

\BookQuestion{网友：您在投资过程中都犯过哪些错误呢？是怎样发现并纠正的呢？}

前面说过一些了，这里补充一个没说过的。去年买过一只ETF（Exchange- Traded Fund，交易所交易基金），叫UNG，是和天然气挂钩的指数。我买的时候天然气的价格只有3块左右，而天然气的长期成本要大约6块以上。我认为没什么东西价格可以长期在成本价以下，就跑去买了一大堆UNG。后来经过认真研究，发现UNG并不能和天然气线性相关，而且时间损耗还很大，就觉得自己错了。但还有点不舍得卖，因为已经亏钱了。再后来想起知错就改，不管多大的代价都是最小的代价的原则，最后还是下决心亏钱了结了。

呵呵，这单我管的所有账号加起来亏了超过 \(\times\)\(\times\) 万，但如果当时没下决心卖的话，到目前为止要大概多亏原来亏损的接近4倍。而天然气现在的价格比我买UNG时要高50\%。这个错误主要是因为对投资标的的了解不够就下手了（当时整体获利情况很好，手里有很多现金）。以后一定不能大意啊。（2010-03-22）

% 来源：EPUB part051.xhtml
\subsection{投资是个概率事件}

我认为“懂”不是个绝对的概念，大概就是“毛估估懂”的意思。懂不懂只有自己才知道，任何别人都很难评价。比如说我觉得我懂通用电气（GE）的文化的价值，但不懂的人就会在其很便宜的时候把股票卖给我。

从投资的角度而言，当你觉得你买的股票掉而你真的不受影响的时候，你大概就是懂了（不是心理因素），不然就有可能是在投机（尤其是当你觉得害怕的时候）。

比如，雅虎（Yahoo）最近从今年的高点掉下来不少，我也重新审视过我投雅虎的理由，发现好的理由并没有任何减少，负面的东西也没有任何增加，我也就没有任何事情需要担心的了。（2010-05-20）

“懂”的定义不要太理想化。“懂”并不是能看到未来一切的“水晶球”，能看到一些重要的东西就很不错了。长远而言，我觉得投资实际上是个概率事件，真正“懂”的人犯错误的概率低，最后的回报高。很多时候不能简单地用某人对某只股票的行为来判断他对投资的理解。（2010-08-30）

所谓投资厉害不厉害实际上是很难定义的，这又不是比赛。

人们对投资“厉害”与否往往喜欢用某段时间的回报的比例来衡量，但这种衡量往往是不科学的，因为有很多因素会影响到短期的回报。下面举几个可能影响短期投资回报的例子。

能力圈：懂投资的人的回报未必能一直比市场好，尤其是有时候能力圈内没得可投时。

margin：有些用margin甚至是用很多margin的人是有可能在赌对的情况下在某个时段表现很好的——但出来混是一定会还的。

把握性：巴菲特说，他的投资一般是在他认为有95\%的把握时才出手（我觉得也许没那么高），但很多人实际上在有60\%的把握时就出手了（或许更低，所以会看到他们深一脚浅一脚的），可是10000个60\%出手的人里面，总能在某个时段找到运气好而回报特别好的，但未必就是因为这些人厉害。

用功程度：所谓懂投资的真正含义，我个人的理解是懂得不投不懂。但投资回报的好坏其实是取决于投资人能找到自己懂的、好的投资标的，这就需要很用功才能经常找到好的目标（巴菲特的标准是一年一个）。像我这种等待苹果掉头上型的用功程度，四年能找到一个就很不错了。（2012-08-27）

\BookQuestion{网友：从商业模式、企业文化、市场价格这三个维度去考察，懂的边界或标准是什么？考察到什么程度才算真正懂？往往会出现您常说的自认为懂其实还不懂的状况，也会出现因欠缺考虑某个关键因子而投资失误的情况。我记得芒格先生有个筛选公司的过滤清单。}

这里没有一个公式，但你真的懂的时候你会知道的。最重要的东西是你能真正忽略市场的波动，聚焦在生意上。想象你打算买一家公司，拿着10年都不上市你依然能睡好觉，那你大概就懂了。或者倒过来也一样，当你觉得一家公司10年内肯定会有大麻烦的时候，你大概不会想买的吧？（2019-05-04）

看懂了一家公司的最简单的标准就是你不会想去问别人“我是不是看懂了这家公司”。（2019-07-24）

没看懂的公司比较容易鉴别，就是股价一掉你就想卖，涨一点点你也想卖的那种。看懂大概就是怎么涨你都不想卖，大掉时你会全力再买进。

看懂并买进的那种公司最好的状态就是不想卖。（2019-04-06）

你只要一直本着不懂不做，你就会少犯错误，时间长了你的结果自然比你不懂也做要好很多，这是要做对的事情的范畴。

如何搞懂是如何把事情做对的范畴，如果实在是搞不懂，就什么都别碰，把钱存银行，表现可能比大部分人强。（2022-02-18）

在把事情做对的过程中，错误是不可避免的，是个学习的过程。每个人都有自己的学习曲线。（2020-12-13）

\BookQuestion{网友：大道觉得什么时候可以抄作业呢？还是小仓位玩玩纯属为了跟踪一下？}

我只偶尔抄一下巴菲特的作业，纯属好玩。不完全懂是不应该赌身家的。不敢赌身家的“投资”其实都很难算得上正经“投资”。（2023-03-12）

如果不懂，作业是抄不了的。如果懂了，不需要抄。投资非常简单，就是买生意。搞懂投资唯一的办法就是搞懂生意。当然，对搞不懂生意的人们理财是另外一个办法，比如买标普500指数或者QQQ这类指数。买基金其实比买公司还难，因为看懂一个公司比看懂某个人要容易很多，而且基金还要收费。（2024-08-08）

\BookQuestion{网友：我花了两年时间看汽车行业，做不到真懂，是不是该反省并放弃？还是有方法可以优化？}

我自己能看懂的公司也是非常少的，看不懂当然只能放弃了。个人认为，看懂一家公司不会比读一个本科更容易。不过，花很多时间去看那些看不懂的公司是不合算的。我在巴菲特那里学到的非常重要的一点就是先看商业模式，除非你喜欢这家公司的商业模式，不然就不要再往下看了，这样能省很多时间。你的汽车行业的例子我其实也不懂，虽然在特斯拉上赚过不少钱，但最后也因为看不懂而全部放弃了。（2019-03-13）

\BookQuestion{网友：真的不懂相对容易分辨，但就是这似懂非懂害死人。}

多数人的问题不在这里。另外，懂的意思差别也很大。很多人一辈子都努力在懂市场，我只是想办法搞懂某些个企业，难度差别非常大。搞懂市场是不可能的任务，Mission impossible！懂了吗？

懂市场这个说法不严谨。市场长期是“称重机”，短期是“投票器”。搞懂这两点的区别很重要。多数人花太多时间和心思在搞懂“投票器”上了。（2022-03-05）

% 来源：EPUB part052.xhtml
\subsection{你懂得，才是你的好生意}

\BookQuestion{网友：大家都知道巴菲特不做科技股，那不符合他的理念。但科技股反而成就了您，您怎样衡量这种反差？}

我只是做我认为我能懂的东西（以为自己懂也不一定就真懂了），有些可能也许正好是大家说的所谓科技股吧。我分不清什么是科技股。（2010-05-30）

看得懂什么不重要。好像现在的中国首富是能看懂饮料的人，也许下一个是个能看懂猪饲料的人。

另外，我从来不认为自己介入过互联网或高科技什么的。互联网或所谓的高科技都只是工具而已，就像当年的铁路、公路、高速公路、航空等等一样。（2010-10-22）

\BookQuestion{网友：您的投资公司名字可以透露吗？想跟踪它的13- F。}

直接问买啥可以马上赚大钱岂不是更简单？知道别人投什么对不懂公司的人来说不会有啥帮助的。有时间还是看看巴菲特1998年在佛罗里达的演讲吧，有视频的。（2024-06-08）

你懂得，所以才是你的好生意。（2010-10-23）

\BookQuestion{网友：巴菲特不喜欢高科技行业是因为变化太快吗？}

巴菲特不喜欢高科技？巴菲特从来没说不喜欢，只说看不懂而已，但他一旦看懂也是一样会出手的，比如IBM。巴菲特还说过，如果你能看懂变化，你将能赚大钱。投资最重要的不是巴菲特能看懂什么，最重要的是你自己能看懂什么。（2013-05-26）

\BookQuestion{网友：低市盈率是价值投资的一种派别，相对来说比较低级。价值投资的高级形式就是对公司进行估值，用4角钱价格买1元钱的价值。需要你对公司真正了解，才能很好地做出正确的评估。}

个人认为投资没有形式上的高低，只有懂的程度的差异。如果我能预见到一个公司有高成长，我当然愿意投。我从来都很愿意投高成长的公司——如果我能确认他会是高成长的。（2010-04-28）

\BookQuestion{网友：投资股票，你最热衷什么行业？什么行业确定性最大？巴菲特提出最好的企业是不需要投入大量资本于更新而能维持优势、创造利润的企业。从这点，很多技术型企业就被踢出去了。而您好像投资了不少技术企业？}

呵呵，行业问题说不清，好像存在的大多行业里都有好公司。

巴菲特有几条东西：能力圈、护城河、安全边际，他没说哪个行业。巴菲特追求的是产品很难发生变化的公司，所以他买了后就可以长久持有。但他也说过，如果你能看懂变化，你将会赚到大钱。归根到底，买股票就是买公司。无论你看懂的是长久还是变化，只要是真懂，便宜时就是好机会。

如果一两年能看到一个这样的好机会，你的回报一定会很好。（2010-05-08）

\BookQuestion{网友：银行股的估值便宜多了，段总就是不相信银行股……呵呵，我对银行没恐惧感，但确实看不懂。}

顺便说下，投资很有趣的地方就是，你能看懂的东西就已经能让你足够忙和得到足够回报了。另外，要想搞懂自己不擅长的东西往往没有那么容易，同样的机会成本（时间）在自己明白的圈子里往往有大得多的回报。（2011-03-04）

\BookQuestion{网友：能讲下小资金投资的方向吗？}

选择你能搞懂的好公司，和资金大小其实无关的。（2019-03-15）

\BookQuestion{网友：能否说说460美元的苹果与160元人民币的茅台哪个更便宜？}

取决于你懂哪个以及懂多少。（2013-03-23）

我能力圈小，能看懂的东西太少了，这些年确实没发现什么能让我用苹果或茅台去换的公司。（2017-03-13）

\BookQuestion{网友：您曾说看好腾讯的商业模式和企业文化，它是家好公司。可是后来您又说，看不懂它未来的现金流。而且您的持仓很低只有1\% （仿佛后来又不是太看好）。我想不明白您看不懂它未来现金流具体是指哪方面，对它不确定的又是什么？}

看不懂未来现金流的意思是看不懂公司未来到底能赚多少钱。投资是要花钱的，所以你需要知道你投资的公司未来能赚多少钱，是不是能赚的比你投的多，不然就不应该投。大部分公司我都看不懂，腾讯我有点懂但不是那么透。

\BookQuestion{网友：按照这个懂的标准，普通人或许很久（3年、5年）也找不到一个生意模式好、又几乎全懂、价格又合适的公司了。}

这也是为什么大道一直觉得自己是普通人的原因。大道过去十几年实际上就真的看懂了两家公司，苹果和茅台。平均六年看懂一家。（2023-11-16）

\BookQuestion{网友：请问您觉得您和李录与巴菲特和芒格在做价值投资上最大的差距在哪？您和李录在哪方面还需要学习？}

最大的差距大概是他们的英语都比我好很多。当然，英语上李录和那两位比（我猜）也会有很大差距的（母语和非母语的差别），但他的差距对他的投资生涯影响不大，和我的差距有本质的差别。

从价值投资的角度看，也许你说的差距用“不同”来说更好些。

从对价值投资的理解上，我认为我们是完全一样的，因为价值投资的本质就一点：买股票就是买公司，买公司就是买公司未来现金流的折现。所有所有所有其他的关于价值投资的说法，其实都是在讨论如何来确定这个未来的现金流折现到底大概会有多少的问题。

我们的不同点主要会体现在能力圈的不同，每个人懂的东西都不同。当然，能力圈的差距也是有的，比如他们做股票投资的时间都比我长很多。

能力圈的意思就是你能够判断未来现金流折现的范围。知道自己的能力圈有多大比能力圈本身有多大要重要得多。

不过，我现在管的钱的绝对值比巴菲特和芒格在我这个年龄时还要多呢，呵呵。不能和李录比，他比我年轻很多啊。

我和他们三个人的最大的不同确实有一点，那就是他们都是职业投资家，我是业余爱好者，对投资的时间安排比例会有点不同。

我们各方面需要学习的东西应该都很多，不过李录大概不需要像我一样还要花时间学英语了。（2011-01-04）

学巴菲特最重要的是学他的原则，毕竟每个人的能力圈不同。就像我敢重仓苹果，巴菲特大概不会一样。我有个朋友在比亚迪上亏了很多钱，我问他为什么买比亚迪，他说巴菲特都买了，这种学法是会出事的。我没买过比亚迪，虽然一直有兴趣看，但一直没看出买的兴趣来，所以巴菲特买不买对我完全没影响。

总之，崇拜任何人都是不对的，最重要的是芒格讲的那个“理性”。（2012-01-09）

举个例子也许你就明白了。巴菲特有很多保险和金融的投资，我基本没有，因为我还不懂，总觉得不踏实。我投了一些和互联网相关的公司，巴菲特没投过，因为他不懂。他认为可口可乐是人们必喝的，我认为游戏是人们必玩的。（2010-04-30）

投资中什么最重要？我个人的理解是缺什么，什么重要。投资最重要的是投在你真正懂的东西上。这句话的潜台词是投在你真正认为会赚钱的地方（公司）。

我对所谓赚钱的定义是：回报比长期无风险债券高。一个人是否了解一个公司能否赚钱，和他的学历并没有必然的关系。虽然学历高的人一般学习的能力会强些，但学校并不教如何投资，因为真正懂投资的都很难在学校任教，不然投资大师就该是些教授了。不过在学校里可以学到很多最基本的东西，比如如何做财务分析等等，这些对了解投资目标会很有帮助。

无论学历高低，一个人总会懂些什么，而你懂的东西可能有一天会让你发现机会。我自己抓住的机会也好像和学历没什么必然的联系。

比如我们能在网易上赚到100多倍是因为我在做小霸王时就有了很多对游戏的理解，这种理解学校是不会教的，书上也没有，财报里也看不出来。我也曾试图告诉别人我的理解，结果发现好难。又比如我当时敢重手买GE，是因为作为企业经营者，我们跟踪GE的企业文化很多年，我从心底认为GE是家伟大的公司。

我说的“任何人都可以从事投资”的意思是我认为并没有一个“只有‘某种人’才可以投资”的定义。但适合投资的人的比例应该是很小的。可能是因为投资的原则太简单，而简单的东西往往是最难的吧。

顺便说一句什么是简单的投资原则：当你在买一只股票时，你就是在买这家公司！简单吗？难吗？（2010-02-06）

\BookQuestion{网友：想起了彼得·林奇说过的话，你喜欢公司的产品，就要想想要不要买它的股票。}

我一般倒过来想，不喜欢的产品的公司的股票不碰，因为没办法了解（感觉）。（2011-03-24）

\BookQuestion{网友：看懂一家公司是否一定要到上市公司所在地考察？有些人为了了解上市公司，就专程到上市公司的所在地租房住了半年。}

公司内部人也未必就懂自己公司，在当地住上半年怎么就一定能懂？搞懂一个公司实际上是没有公式的，不然就可以在学校开课教了。（2011-01-13）

过多拜访公司容易短视。我的意思是，很多人拜访公司是为了了解公司短期运营的状况，这种心态对长期投资的帮助不一定大。（2010-04-12）

“永远不要问理发师你是否需要理发。”（巴菲特）

知道巴菲特为什么认为拜访公司不一定有用了吧？（2012-07-27）

所谓价值投资说通了就是要能理解企业。所有能帮助你理解企业的东西都是有用的，但不会有一个充分条件保证你会投资。我没啥更好的建议。我一般讲的都是要有自己不做的事情，至于怎么才能把事情做对是要靠自己去摸索的。（2019-06-26）

\BookQuestion{网友：我有很多东西不懂。研究一家公司，通过看年报、公司官网，读企业传记和亲身体验公司产品等方式了解到的有关企业文化与商业模式的东西总觉得太少了，也看不懂以后的现金流（看不出公司在未来的竞争中处于一种什么状态）。请教大道还有什么方法可以进一步加深研究吗？}

其实是有的，大概就是先开一家公司，干个几年后就明白了。千万不要以为这是玩笑。黄峥当年问过我类似的问题，就是不管怎么看公司总觉得差点啥，当时我们说的就是经营一家公司可能会有帮助。后来他就回国开了公司了。（2020-12-03）

\BookQuestion{网友：如何知道自己看懂（或者没有看懂）某个企业呢？比如茅台，记得你说过自己并不喝白酒，也没有听说过你曾开过酒厂。}

一般情况下我也是通过公开渠道了解的。现在网络这么发达，几乎所有东西都能找到，但自己需要会辨别。比如苹果，你可以看下苹果所有的产品发布会，你再对照一下同期别的同行的你就会知道区别。企业主页也是一个了解的途径。新闻要小心，因为作者往往带有观点，如果你不明白，很可能会把你带沟里去。功夫花到家了你就会明白的。

从另一个角度再回答一下你的这个问题，我看公司主要看两个东西，生意模式和企业文化。这两样东西中任意一样我不喜欢我就不会再继续看下去，所以没有懂不懂的问题，我不需要懂我没兴趣的公司。如果生意模式我喜欢（当然至少要懂了的意思），企业文化也很好，那就老老实实等好点的价钱。持有一个自己懂且喜欢的公司，最大的特征就是可以完全漠视市场的变化，就是可以拿着睡好觉。买个“票”然后天天睡不好觉，到处问别人怎么看，每天网上到处找各种相关消息的那种就是不懂。搞懂一家好公司其实是非常不容易的，多数情况下不是简单看下报表就可以的，但偶尔确实有那种简单看下报表就捡到宝的情况，比如巴菲特的某些公司，但那其实也是很多年碰上一个的概率，可遇不可求的。（2019-04-07）

\BookQuestion{网友：投资难在未来风险的不确定性。}

不确定性是用安全边际去对付的。（2011-05-09）

\BookQuestion{网友：我认为安全边际是重中之重，您怎么看？}

很重要的东西之一吧。（2010-04-04）

我对安全边际的理解是对公司的理解度而不是股票短期的价格波动。对一个伟大公司而言，有时候一点点价格差别10年后看都不是事，但因为差那点价格而错过一个好公司就可能是大事了。

其实这个问题我想了很多年，一直没懂巴菲特说的安全边际是什么意思，但自己最后的体会就是这样，不了解的公司再便宜也可能不便宜。（2020-10-15）

“如果你对一家企业及其未来发展了如指掌，那么你所需要的安全边际就会很小。所以，假设你仍想投资一家较为脆弱的企业，你就需要更大的安全边际。就好比你开着一辆卡车通过一座标明限重10000磅的桥，而你的卡车重9800磅，如果桥下只有6英寸深的沟壑，你可能会觉得没问题；但要是这座桥横跨大峡谷，你可能就会觉得需要更大的安全边际。”（巴菲特，1997）

以前没看过巴菲特讲的原话，今天偶尔看到，发现跟我的意思一样一样的。（2024-01-21）

% 来源：EPUB part053.xhtml
\subsection{高手在于错误率低}

“今天重新看了一遍《赢在中国》，听到了马云的一句话很矛盾，他说不犯错将是最大的错，而巴菲特先生的原则却是不要犯错，这要怎么理解怎么把握呢？”这段话引自网友，我觉得是个很好的问题，故说说自己的看法。

我不知道这里所提巴菲特的原话是什么，我只记得巴菲特说过大概这样的话：第一，不要亏钱；第二，不要忘记第一条。我个人理解他老人家在这里讲的是投资安全的问题。如果一个人投资到他不懂的东西里就有可能亏钱，所以要避免。这并不意味着巴菲特就不会亏钱，事实上他亏过的钱的总额比我们谁都多，因为他也有犯错误的时候。呵呵，光是他投在康菲石油上亏的钱就有几十亿美元啊。我一直没看明白他为什么要投康菲石油，也许他也犯了以为看懂实际却没看懂的错误？但是，巴菲特有今天，我个人觉得最重要的就是他很少犯“原则性错误”，也就是说，他认为是他能力以外的事他坚决不碰。几十年下来，他犯的错远远少于同行，仅此而已。

巴菲特之所以是巴菲特，最重要的就是他能坚持做正确的事，也就是原则性错误的事不做。比如，他认为自己不懂的事就坚决不做，我从巴菲特那里学到的最重要的东西也就是这个。

马云讲的是另外一面，是如何把事情做对的问题。无论是谁，在把事情做对的过程当中都是有可能犯错的，好像也没听说哪个所谓成功人士没犯过错。把事情做对的过程往往是一个学习的过程，犯错往往是不可避免的。不犯错的办法只有一个，那就是什么都不做。什么都不做有时可能就是最大的错。

不知道这么解释你是否能明白？有时想想中文真是很有意思。一会儿说“慈不带兵”，一会儿说“爱兵如子”；一会说要“三思而后行”，一会又说要“当机立断”。呵呵，好像并不矛盾啊。他们说的都是事物的不同面而已吧？（2010-03-14）

投资的一般规律是：出手越多，赚得越少，或赔得越多。（2010-04-23）

\BookQuestion{网友：您大量买入网易和U- Haul Holding Company（UHAL）似乎都发生在2002、2003年左右，最近三五年有没有比较经典的大手笔投资的案例（除了苹果）？}

最好把网易也除掉，我就快成笑话了。（仿芒格语）（2012-02-02）

平均两年做一个大的投资决策的效果肯定比一年做两个更好。（2012-06-26）

大道这四年里只真正出手了一次，就是苹果，上一次出手是GE，大概这是其会成为极少数人的原因？

如果大道只敢平均两年出手一次，凭什么“散户”可以出手更多？（2012-07-17）

我选股不厉害，是业余级的，花的功夫太少。但我犯错的机会的确是大师级的，不懂就是不碰。真看得懂这两句话的人日子长了一定会有不错的赚头哈。（2011-11-03）

好的价值投资者并不是那些没有犯过错误的人。好的价值投资者的最典型特征是从很多年的历史来看，犯错误（尤其是大错）的概率很低。常胜将军的意思是赢的比例高而已，百战百胜则是神话。（2011-12-19）

经营企业和投资一样，少犯错很重要。但所谓少犯错不是通过什么都不敢做实现的，那叫裹足不前。少犯错要通过坚持做对的事情来实现。而所谓坚持做对的事情是发现是错的事情就马上停止，不管多大的代价都会是最小的代价。不为清单也很重要。（2020-10-14）

\BookQuestion{网友：看伯克希尔2002年的年报，巴菲特说：“在投资股票时，我们预期每一笔投资都会成功，在我们经营伯克希尔的38年中，投资获利的个案和投资亏损的个案比例约为100：1。”这才是巴菲特真正牛的地方啊！厉害在几乎不犯错。}

呵呵，你说到关键了。高手和其他选手的差别就在错误率低而不是能打出多少好球来。（2011-09-07）

\BookQuestion{网友：比尔·米勒（Bill Miller）这次金融危机亏掉了老本，连续15年战胜指数，但一年的亏损就把15年的利润搞没有了，亏的比赚的多。他是不是太高估自己了？以前觉得他是个价值投资者，这次危机中露出原形了，他曾经在安然公司上栽跟头，这次又在房地美等企业上栽跟头。其实价值骗子还蛮多的。所以你对多少次都无所谓，关键是错的时候千万不要损失太大，这涉及安全边际。}

我认为当一个人开始追求战胜指数的时候，他可能已经失去平常心了。很多人在这次金融危机中都有一些再也回不来的重伤股，巴菲特一个也没有，这绝对不是偶然的，但好像很少人注意到这点。（2010-05-23）

\BookQuestion{网友：随着投资经验的增加，你和以前比进步在哪些方面？}

呵呵，心态好很多，不怕失去机会，最重要的是不犯大错。（2010-03-18）

\BookQuestion{网友：投资和打好高尔夫，哪个更难？}

两个都是我最喜欢的游戏，它们都很难。也许高尔夫更难些吧。我认为高尔夫更难的原因是，高尔夫你不能说这杆我还没想好，所以我现在不想打。投资你永远都可以说，我还没想好，我不想做什么。（2019-03-30）

% 来源：EPUB part054.xhtml
\subsection{做空是愚蠢的}

其实我最大的财富就是我犯过的这些错误，巴菲特也是。（2010-05-05）

\BookQuestion{网友：能否分享一下自己认为的最差的一次投资？}

呵呵，亏钱的投资有好几次了，多数都可以原谅，但有一次属于极度愚蠢，就是跑去做空百度。之前投资表现非常好，有点飘飘然，真以为自己很厉害。开始还是想小玩玩，后来又不服输，最后被夹空而投降，所有账号加起来亏了很多钱，大概有1.5-2亿美元，其中有个账号到现在都翻不了身。

最可惜的是，这次错误把我们的现金储备全部消耗掉了，机会成本巨大，不然这次金融危机我就可以帮大家赚更多钱了。而且最遗憾的是，这个错误是发生在巴菲特叫我不要做空之后。现在知道什么叫不听老人言的后果了吧？值得欣慰的是，这三年总的表现还不错。真有点感谢金融危机给了我逆转的机会。

巴菲特的教导千万别忘了：不做空，不借钱，不做不懂的东西！

不过，这其实不算我最差的投资，因为这根本就是彻头彻尾的投机。以后不会再有这种例子了。（2010-03-07）

\BookQuestion{网友：“最可惜的是，这次错误把现金储备全部消耗掉了，机会成本巨大”，是不是指现金用完了，真正好的投资标的来时，却没钱买了。}

是的，机会成本大，亏点钱是小事。（2010-05-12）

\BookQuestion{网友：当时卖空百度的理由是什么？}

首先，卖空行为是错的。当时确实看到了一些百度不对的东西，事实上百度后来也因为这个掉得很低（如果我能一直空到那时实际上可以赚大钱），但我也确实忽略了一些对百度有利的重要条件，比如政策环境等。总而言之，卖空行为是投机行为，吾不该为之。（2010-09-14）

当时不知道哪根神经起了主导作用，属于失去平常心的愚蠢行为。（2012-03-14）

\BookQuestion{网友：基于对价值的估计，在价格过高时做空算不算价值投资？}

做空不是价值投资，因为你经常需要面对市场的疯狂，最糟糕的是，你不知道市场到底有多疯狂，所以做空会睡不好觉的。总而言之，价值投资就是那种能睡好觉的投资。（2011-05-22）

任何时候，只要你还想着要空一把谁，那就表示你还是个投机分子（以为自己比市场聪明）。早点收起做空头的念头吧，老老实实去找那些可以让你睡得好觉的deal！（2011-10-29）

做空是愚蠢的！（2013-03-04）

我不做空，不想跟自己过不去。做空一家公司需要了解的东西远比做多要难得多得多，而且一旦错了就可能是灾难性的。（2020-10-28）

希望你有一天能理解巴菲特为什么讲不要做空，那时，你的投资就会再进一步。（2010-05-22）

我有个邻居（没那么熟），是专门做做空生意的（我搞不懂怎么做哈）。前段时间突然问我是否知道“跟谁学”这家公司，说好几个做空公司在做空它，因为它做假账，理由是他们的成长太快了。我真的让我们教育电子的同事去了解了一下，得出的结论是：它的成长在这个环境下看不出任何明显不符合逻辑的地方。于是我告诫邻居，中国很大，这点营业额的成长比例其实不是不可能的。那时候股价好像是在30-40之间。前两三个礼拜打球又碰到邻居，问起他是否还空着呢？他说是的，他们的成长太不可思议了，肯定是做假账。我看了一眼股价，那天80左右。

我完全不了解跟谁学这家公司，这里我想说的是：最好别做空哈！这个世界有很多你不懂的事情，为什么要跟自己过不去？

我其实就是想告诉那些和我的邻居有同样想法的人们，做空是很危险的，时间很可能不站在你这边。我看到这家公司从30多涨到130多的时候，我真的有点替邻居担心哈，我觉得他恐怕很难睡好觉，何苦呢？！

不喜欢的公司避开就好了，千万不要去做空，因为看错一次会让你难受很久甚至一辈子的，尤其是那些已经有钱的人们。没人希望自己需要富两次的。（2020-08-11）

% 来源：EPUB part055.xhtml
\subsection{不用margin是投资的基本要求}

如果你懂投资，你不需要用杠杆，因为你早晚会变富的。如果你不懂投资，你更不应该用杠杆，不然裸奔的将可能会是你。投资是件快乐的事情，用杠杆会让你有机会睡不好觉的。（2020-12-06）

\BookQuestion{网友：做生意可以贷款，为什么不建议借钱买股票？}

借钱做生意一般而言是有还款期的，而且借款人对要做的生意一般了解度高，但也还是有很多人因为借钱做生意而陷入一生的麻烦。

借钱做股票是用你的其他股票抵押的，完全由经纪人（broker）控制，但市场大跌时他们会随时要求你卖股票还钱。股票市场有时会疯狂到无法想象，用margin的人一生也许只需要碰上一次市场向下的疯狂就可能会损失惨重（wiped out）。芒格这么聪明的人就曾经（20世纪70年代）因为margin而陷入大麻烦。巴菲特厉害的地方是其投资生涯中从来就没有过致命的大麻烦。巴菲特讲过，即使只有1\%的破产机会你也不应该干。懂投资的人不需要用margin，不懂投资的人更不应该用margin。（2011-09-08）

\BookQuestion{网友：决定申请融资融券账户了！段大哥一直说不要借钱，还是忍不住啊！}

从此你会有个激动人心的生活！（2012-08-10）

\BookQuestion{网友：如果20年前融资买的是苹果、茅台、腾讯这些好公司，好像也没问题。但加杠杆的人总是喜欢乱买。}

苹果在过去二十年里，从高点下降40\%或以上就有至少10次， 10\( - \)20\%更是不计其数，要是用margin买了苹果，就算不破产，大概率也已经吓破胆了？（2023-12-19）

\BookQuestion{网友：我长期低杠杆投资比如0.3倍以内，融资成本6个点，在股指偏低位建仓，国内融资盘大概要跌七成才会爆仓，这种低杠杆长线价值投资策略可以吗？}

勿以恶小而为之。（2019-05-20）

用margin的人都不傻，但用margin是会上瘾的，直到掉坑里为止。其实还是巴菲特说得好，如果你懂投资，你不需要用margin，因为你早晚会富有的。（2024-09-03）

长期资本（Long-Term Capital Management）就是个很好的例子。那些人都极度聪明，两个诺贝尔经济学奖获得者加一群特别“厉害的人”，而且还赌上了自己的身家。（2024-09-01）

\BookQuestion{网友：2023年时等了大半年茅台就是不跌，忍不住1600元时用手上的钱都买了。最近倒是一路跌下来了，可是没有余钱加仓了，好想借钱抄底啊。}

我们不知道市场会有多疯狂，也不知道事情可能会坏到什么程度。不要用margin是投资的基本要求！我们一直说的是10年20年，但你用margin可能会让你在那之前就陷入困境，没必要。（2024-09-18）

\BookQuestion{网友：A股大动荡，现在多庆幸自己没融资，压力小多了。真到了危机时刻，才会发现您说的原则是多么珍贵。}

把事情做对的能力只有在做对的事情时才能让时间站在自己的一边，不然早晚会出事的。（2015-08-10）

反正你借不借钱一生当中都会失去无穷机会的，但借钱可能会让你再也没机会了。（2010-04-11）

用杠杆炒房和用margin没有本质差别，但问题来的时候要更惨，因为房子的流通性一定会差些。（2013-10-09）

\BookQuestion{网友：感觉很惭愧，这些道理都是从十几年前就看熟了的，能背的，竟然还会时不时“不小心”动了用margin的念头，好在没有真的去做。}

你值得惭愧一下。理解赚慢钱不容易，但理解复利更难。慢钱在复利的影响下，效果是非常厉害的。（2024-09-04）

\BookQuestion{网友：如果投了一个负债率很高的企业，属于使用margin吗？}

似乎不属于直接用margin，不会导致你需要富两次。但是，我是不太喜欢买负债率高的企业的，除非我能特别懂。（2016-11-12）

“如果你想提高自己的认知能力，忘记过去犯过的错误是坚决不行的。”（芒格）

想起小偷被抓以后，满脑子都是如何改进偷技的想法的那个故事。

所以这里的错误有两种性质：（1）做了错的事情——以后不能再做——放上不为清单；（2）在把事情做对上犯了错误——这个不可避免，但可以通过学习改进。那个小偷犯的是第一个错，但用的是第二个改进办法，会更糟的。

“最好是从别人的悲惨经历中学到深刻教训，而不是自己的。我们有些成功是早就预言的，有些是意外获得的。”（芒格）

不听老人言，吃苦在眼前和眼后。看了这句话以后，其实大多数人还是不会听的。

能从自己的错误里吸取教训就已经很了不起了，能从别人的错误里吸取教训的那都是天才。

坚持犯自己犯过的错误的人也不用太不高兴啊，至少股市上85\%以上的人是这样的。有人陪总是一件不错的事情。（2012-06-26）
